% SEAD HTN decomposition (illustrates Table tab:htn-trace)
% include with % SEAD HTN decomposition (illustrates Table tab:htn-trace)
% include with % SEAD HTN decomposition (illustrates Table tab:htn-trace)
% include with % SEAD HTN decomposition (illustrates Table tab:htn-trace)
% include with \input{chapters/htn_trace_tree}
% Uses the compound/primitive/method styles and \htnpre from main.tex
\begin{figure}[htbp]
  \centering
  \tikzset{
    stepno/.style   ={circle, fill=green!45!black, text=white, font=\scriptsize\bfseries,
                    inner sep=0pt, minimum size=12pt},
    notused/.style={draw=black!30, text=black!40, fill=black!3},
  }
  \begin{forest}
    for tree={
      font=\footnotesize, align=center, inner sep=3pt, rounded corners=2pt,
      minimum height=0.7cm,
      l sep=0.8cm, s sep=0.6cm,
      edge={draw=green!45!black, very thick},
      parent anchor=south, child anchor=north,
    }
    [{\textbf{suppress-enemy-radar}}, compound, label={[stepno]right:1}
      [{m0: standard SEAD\htnpre{drone ready}}, method
        [{fly-to-\\area}, primitive, label={[stepno]below:2}]
        [{find-radar}, compound, label={[stepno]right:3}
          [{m1: passive\\search\htnpre{radar emitting}}, method, notused,
            edge={draw=black!30, thick, dashed}
            [{listen-\\passively}, primitive, notused,
              edge={draw=black!30, thick, dashed}]
          ]
          [{m2: decoy\\search\htnpre{radar silent}}, method
            [{act-as-\\decoy}, primitive, label={[stepno]below:4}]
            [{listen-\\passively}, primitive, label={[stepno]below:5}]
          ]
        ]
        [{suppress-\\radar}, primitive, label={[stepno]below:6}]
        [{check-\\result}, primitive, label={[stepno]below:7}]
        [{return-to-\\base}, primitive, label={[stepno]below:8}]
      ]
    ]
  \end{forest}

  \medskip
  {\small $s_0 = \{\htnfact{drone-ready},\ \htnfact{at-base},\ \htnfact{radar-silent}\}$}

  \caption[Decomposition of the SEAD HTN]{%
    Decomposition of the SEAD HTN for the initial state $s_0$.
    Highlighted edges show the chosen decomposition; method \textit{m1}
    is not applicable because the radar is silent. Numbers refer to the
    steps in Table~\ref{tab:htn-trace}.}
  \label{fig:htn-trace-tree}
\end{figure}

% Uses the compound/primitive/method styles and \htnpre from main.tex
\begin{figure}[htbp]
  \centering
  \tikzset{
    stepno/.style   ={circle, fill=green!45!black, text=white, font=\scriptsize\bfseries,
                    inner sep=0pt, minimum size=12pt},
    notused/.style={draw=black!30, text=black!40, fill=black!3},
  }
  \begin{forest}
    for tree={
      font=\footnotesize, align=center, inner sep=3pt, rounded corners=2pt,
      minimum height=0.7cm,
      l sep=0.8cm, s sep=0.6cm,
      edge={draw=green!45!black, very thick},
      parent anchor=south, child anchor=north,
    }
    [{\textbf{suppress-enemy-radar}}, compound, label={[stepno]right:1}
      [{m0: standard SEAD\htnpre{drone ready}}, method
        [{fly-to-\\area}, primitive, label={[stepno]below:2}]
        [{find-radar}, compound, label={[stepno]right:3}
          [{m1: passive\\search\htnpre{radar emitting}}, method, notused,
            edge={draw=black!30, thick, dashed}
            [{listen-\\passively}, primitive, notused,
              edge={draw=black!30, thick, dashed}]
          ]
          [{m2: decoy\\search\htnpre{radar silent}}, method
            [{act-as-\\decoy}, primitive, label={[stepno]below:4}]
            [{listen-\\passively}, primitive, label={[stepno]below:5}]
          ]
        ]
        [{suppress-\\radar}, primitive, label={[stepno]below:6}]
        [{check-\\result}, primitive, label={[stepno]below:7}]
        [{return-to-\\base}, primitive, label={[stepno]below:8}]
      ]
    ]
  \end{forest}

  \medskip
  {\small $s_0 = \{\htnfact{drone-ready},\ \htnfact{at-base},\ \htnfact{radar-silent}\}$}

  \caption[Decomposition of the SEAD HTN]{%
    Decomposition of the SEAD HTN for the initial state $s_0$.
    Highlighted edges show the chosen decomposition; method \textit{m1}
    is not applicable because the radar is silent. Numbers refer to the
    steps in Table~\ref{tab:htn-trace}.}
  \label{fig:htn-trace-tree}
\end{figure}

% Uses the compound/primitive/method styles and \htnpre from main.tex
\begin{figure}[htbp]
  \centering
  \tikzset{
    stepno/.style   ={circle, fill=green!45!black, text=white, font=\scriptsize\bfseries,
                    inner sep=0pt, minimum size=12pt},
    notused/.style={draw=black!30, text=black!40, fill=black!3},
  }
  \begin{forest}
    for tree={
      font=\footnotesize, align=center, inner sep=3pt, rounded corners=2pt,
      minimum height=0.7cm,
      l sep=0.8cm, s sep=0.6cm,
      edge={draw=green!45!black, very thick},
      parent anchor=south, child anchor=north,
    }
    [{\textbf{suppress-enemy-radar}}, compound, label={[stepno]right:1}
      [{m0: standard SEAD\htnpre{drone ready}}, method
        [{fly-to-\\area}, primitive, label={[stepno]below:2}]
        [{find-radar}, compound, label={[stepno]right:3}
          [{m1: passive\\search\htnpre{radar emitting}}, method, notused,
            edge={draw=black!30, thick, dashed}
            [{listen-\\passively}, primitive, notused,
              edge={draw=black!30, thick, dashed}]
          ]
          [{m2: decoy\\search\htnpre{radar silent}}, method
            [{act-as-\\decoy}, primitive, label={[stepno]below:4}]
            [{listen-\\passively}, primitive, label={[stepno]below:5}]
          ]
        ]
        [{suppress-\\radar}, primitive, label={[stepno]below:6}]
        [{check-\\result}, primitive, label={[stepno]below:7}]
        [{return-to-\\base}, primitive, label={[stepno]below:8}]
      ]
    ]
  \end{forest}

  \medskip
  {\small $s_0 = \{\htnfact{drone-ready},\ \htnfact{at-base},\ \htnfact{radar-silent}\}$}

  \caption[Decomposition of the SEAD HTN]{%
    Decomposition of the SEAD HTN for the initial state $s_0$.
    Highlighted edges show the chosen decomposition; method \textit{m1}
    is not applicable because the radar is silent. Numbers refer to the
    steps in Table~\ref{tab:htn-trace}.}
  \label{fig:htn-trace-tree}
\end{figure}

% Uses the compound/primitive/method styles and \htnpre from main.tex
\begin{figure}[H]
  \centering
  \tikzset{
    stepno/.style   ={circle, fill=green!45!black, text=white, font=\scriptsize\bfseries,
                    inner sep=0pt, minimum size=12pt},
    notused/.style={draw=black!30, text=black!40, fill=black!3},
  }
  \begin{forest}
    for tree={
      font=\footnotesize, align=center, inner sep=3pt, rounded corners=2pt,
      minimum height=0.7cm,
      l sep=0.8cm, s sep=0.6cm,
      edge={draw=green!45!black, very thick},
      parent anchor=south, child anchor=north,
    }
    [{\textbf{suppress-enemy-radar}}, compound, label={[stepno]right:1}
      [{m0: standard SEAD\htnpre{drone ready}}, method
        [{fly-to-\\area}, primitive, label={[stepno]below:2}]
        [{find-radar}, compound, label={[stepno]right:3}
          [{m1: passive\\search\htnpre{radar emitting}}, method, notused,
            edge={draw=black!30, thick, dashed}
            [{listen-\\passively}, primitive, notused,
              edge={draw=black!30, thick, dashed}]
          ]
          [{m2: decoy\\search\htnpre{radar silent}}, method
            [{act-as-\\decoy}, primitive, label={[stepno]below:4}]
            [{listen-\\passively}, primitive, label={[stepno]below:5}]
          ]
        ]
        [{suppress-\\radar}, primitive, label={[stepno]below:6}]
        [{check-\\result}, primitive, label={[stepno]below:7}]
        [{return-to-\\base}, primitive, label={[stepno]below:8}]
      ]
    ]
  \end{forest}

  \medskip
  {\small $s_0 = \{\htnfact{drone-ready},\ \htnfact{at-base},\ \htnfact{radar-silent}\}$}

  \caption[Decomposition of the SEAD HTN]{%
    Decomposition of the SEAD HTN for the initial state $s_0$.
    Highlighted edges show the chosen decomposition; method \textit{m1}
    is not applicable because the radar is silent. Numbers refer to the
    steps in Table~\ref{tab:htn-trace}.}
  \label{fig:htn-trace-tree}
\end{figure}
